\section{Conclusion and Limitations}
\label{sec:conclusion}

Probing is an intervention: it reveals hidden dynamics while perturbing the state to control. In our Genesis benchmark, \PTA{} improves on recoverable elastoplastic material and pays transfer costs on sand, snow, and shifted sand. A matched-encoder six-seed audit strengthens the elastoplastic direction, while the asymmetry itself maps where probing should be invoked and where its physical cost dominates.

Recoverability predicts three falsifiable tests: matched materials should flip sign by measured relaxation rather than label; stronger probes or shorter settle times should eventually turn elastoplastic gains negative; and dummy probes or non-sand training should not create advantages on persistent materials. Failure of these directional tests would reject recoverability as the organizing explanation.\label{sec:conclusion:end}
